% Section 4, Instances (about 1,000 words; docs/paper-mock.md): the EMBER instance (what is recorded, the three
% repairs, the import-overwrite patch, the edit catalogue). Sources: docs/formalization.md Section 9,
% docs/item-layer-design.md, docs/direction-A-threat-model.md, docs/upstream-issues.md. The requirement checks
% (single-byte flips, random writes) and the population/closure caveats go to the appendix.
% A second instance from another domain, if one is added, gets a short subsection here (about 250 words); which
% domain is interchangeable.

Our main instance is EMBER~\cite{anderson2018ember} in its 2024 version~\cite{joyce2025ember}, the standard static
feature set for Windows executables. Its extractor computes 2,568 features from a PE file: the byte histogram and
byte-entropy histogram, statistics of the printable strings, header fields, per-section statistics, imported
libraries and functions, exports, data directories, and the Authenticode signature. We instrument it without changing
what it computes: each feature group gains a method that returns its items from the same parse, and each feature a
declared readout.

\subsection{What the Instrumented Extractor Records}
\label{sec:instances:record}

Of the 2,568 features, 2,480 are modeled; the residual is the 88 warning flags the PE parser raises, which have no
items. The scanned kinds are bytes, byte-entropy windows (2,048 bytes at steps of 1,024), printable strings, and the
Rich header, which the parser searches for within a fixed range. The parsed kinds are sections, the overlay, header
fields, imports, libraries, exports, the certificate table, and the entry section, which a selection readout picks.
Every readout type of \Cref{tab:fw:readouts} occurs except presence, rate, and standard deviation. On 2,000 files,
the readouts applied to the recorded sums reproduce the original extractor bit for bit on all 2,480 modeled features.

Most of the instrumentation effort goes into locating positions. For each kind, they are the header fields and table
bytes the parser reads to find its items, together with kind-level positions such as the section count and the
bounds of the Rich header's search range. Some are easy to miss: the overlay starts where the last section or data
directory ends, so every section's and directory's extent locates it. Others lie past the end of the file, when a
section table points to raw data that a truncated file no longer contains; on such a file, appending is structural.
We completed the record by changing single bytes and checking whether the parse changed where no locating position
was recorded. A first round found three parser rules the record had missed; after recording them, no flip among
58,666 across headers and directory data changed the parse unannounced. The margins are those of
\Cref{sec:fw:side-effects}: strings reach the nearest non-printable byte, windows reach the windows overlapping the
write (and, under appending, the last 1,024 bytes before the old end), and the Rich header's margin is its whole search
range.

\subsection{Three Repairs}
\label{sec:instances:repairs}

EMBER breaks self-contained contributions (R2) in three places. Two are matters of bookkeeping. Each section's size
is divided by the file's size, so a section's contribution would change whenever the file grows; we let the section
contribute its raw size and move the division into the readout, where the file's size is itself a sum. The imported
libraries are kept as a set; we keep a count per library and read out the number of libraries with a positive count.
Both repairs leave every feature value unchanged.

The third is a bug. EMBER~2024 stores each library's imported functions under the library's name, so when a file has
two import descriptors for the same library, the second replaces the first one's functions. The original EMBER
guarded against this, and the reimplementation lost the guard. The effect is not subtle: adding one descriptor that
imports \texttt{Sleep} from \texttt{KERNEL32.dll} to a file that already imports 114 functions from it lowers the
import count from 116 to 3, and hides the other imports from the model. No model of contributions can repair a
contribution that depends on whether another item exists, so we patch the extractor, as a pure addition that restores
EMBER's original behavior, and run every experiment on the patched variant. On natural files the patch rarely
matters; for an attacker it is a way to remove features, which no add-only threat model would anticipate.

\subsection{Reading the Sums Directly}
\label{sec:instances:sums}

Because the record exposes the sums, a detector need not read the features at all. Most features that can move both
ways do so because of their readout, not their items: a normalized histogram can fall when any bin outside it rises.
Reading the sums instead keeps their directions. We use four such re-encodings: raw byte counts instead of normalized
frequencies; counts over section bodies and the overlay only, the regions no in-place operation writes; signed
section-name buckets split into their positive and negative parts; and import counts with the patch above. Each is
an identity readout over recorded sums, so \Cref{thm:fw:directions} gives its directions with no rule of its own, and
the same bit-exact check covers it.

\subsection{The Edit Catalogue}
\label{sec:instances:catalogue}

The threat model is a catalogue of the twelve operations the public attack toolkits apply to PE
files~\cite{demetrio2021adversarial,song2022mab,kreuk2018deceiving,suciu2019exploring}, with all content (bytes,
section names and properties, imported names) chosen by the attacker within a byte budget
(\Cref{tab:instances:catalogue}). They form three nested threat models: \emph{additions}, which monotone detectors
are designed for; \emph{shifts}, which adds the two insertions that move existing content; and \emph{full}, which adds
every in-place overwrite and removal the toolkits use. Rewriting code and packing are outside all three: they are
whole-file rewrites, undeclared structural edits by nature.

\begin{table}[t]
  \caption{The edit catalogue. Operations stated in elements get their signatures per file from the item table;
  the others declare their item deltas.}
  \label{tab:instances:catalogue}
  \small
  \begin{tabular}{@{}lll@{}}
    \toprule
    Operation & Stated in & Threat model \\
    \midrule
    append bytes at the end & elements & additions \\
    add a section (any name, properties, content) & items & additions \\
    add imports (any library and function names) & items & additions \\
    insert a gap before the PE header & items & shifts \\
    insert a gap before the first section & items & shifts \\
    overwrite section slack & elements & full \\
    overwrite the DOS header and stub & elements & full \\
    overwrite loader-ignored header fields & elements & full \\
    overwrite the checksum & elements & full \\
    rename a section (any 8-byte name) & items & full \\
    remove the debug directory entry & items & full \\
    remove the signature directory entry & items & full \\
    \bottomrule
  \end{tabular}
\end{table}

Each operation stated in items declares its delta, including the fields it must update: adding a section raises the
section count and the image size and writes a new section table entry. Each operation stated in elements declares only
the region it writes and its content domain, and the engine derives the rest. All edits are written by our own
byte-level code, which does exactly what the operation declares. A general editing library such as LIEF~\cite{lief}
rebuilds the file's layout when it adds a section, which breaks honest tools (R5).

Directions hold over a population (\Cref{sec:fw:directions}), and ours is the files whose layout is plain: nothing
the parse locates lies past the end of the file, sections lie within the image, and the headers hold only headers.
Of 2,000 files, 1,674 qualify; most of the rest are files our dataset truncates at 4~MB. A detector checks the
condition at inference and sends the other files to an unconstrained model, which carries no guarantee
(\Cref{sec:evaluation}).
